\subsection{Milestone 6: Monitoring}
\label{sec:m6}
\guidance{
  \textbf{Template:} Describe how the different software engineering practices and tools have been applied.
  (The course template labels this section \enquote{Milestone 5}; it is Milestone~6.)\\
  Monitoring of the running system, data and model (degradation, drift, anomalies) is a typical M6 EDN entry.
}

\subsubsection{Resource monitoring}
\label{sec:m6-resources}
\guidance{
  \textbf{Template:} e.g.\ Grafana, Prometheus.\\
  \textbf{Rubric (5 pts):} To what extent does our monitoring anticipate real-world challenges and provide actionable insights into system performance?\\
  The lab schedule also lists Better Uptime for availability checks.
}

\tbd{resource monitoring setup and dashboards}

\subsubsection{Model performance}
\label{sec:m6-model-performance}
\guidance{
  \textbf{Template:} e.g.\ Alibi Detect.\\
  \textbf{Rubric (5 pts):} To what extent does our monitoring provide timely, actionable signals about data changes, model performance degradation, and user impact?
}

\tbd{drift detection and model performance monitoring}

\subsubsection{Cycle: feedback loops and retraining}
\label{sec:m6-feedback}
\guidance{
  \textbf{Template:} Description of all the stages of the system pipeline, and how they are connected.
  Explain how the pipeline can be used to automate the retraining process.
}

\tbd{end-to-end pipeline and automated retraining}
